\uuid{jgLp}
\titre{Limite}
\auteur{bodin ; hanani ; mzari}
\organisation{exo7}
\licence{CC BY-NC-SA 4.0}
\datecreate{2020-01-20}
\niveau{L1}
\module{Analyse}
\chapitre{Suite}
\sousChapitre{Autre}
\difficulte{1}
\theme{Suites, limite}
\exo7qcmid{lille-1.suites.003}

\contenu{
\begin{question}
\qclassification{121.00}
\qtags{motcle=limite}
Soit $\displaystyle u_n=\frac{n^2+1}{2n^2-1}$ et $\displaystyle v_n=\frac{2n+1}{n^2-1}$. Quelles sont les bonnes réponses ?
\begin{answers}
    \good{$\displaystyle \lim _{n\to +\infty}u_n=\frac{1}{2}$ et $\displaystyle \lim _{n\to +\infty}v_n=0$ }
    \bad{$\displaystyle \lim _{n\to +\infty}u_n=2$ et $\displaystyle \lim _{n\to +\infty}v_n=0$}
    \bad{$\displaystyle \lim _{n\to +\infty}u_n=\frac{1}{2}$ et $\displaystyle \lim _{n\to +\infty}v_n=2$}
    \bad{$\displaystyle \lim _{n\to +\infty}u_n=2$ et $\displaystyle \lim _{n\to +\infty}v_n=+\infty$}
\end{answers}
\begin{explanations}
D'abord, $\displaystyle u_n=\frac{n^2\left(1+\frac{1}{n^2}\right)}{n^2\left(2-\frac{1}{n^2}\right)}=\frac{1+\frac{1}{n^2}}{2-\frac{1}{n^2}}$. Or, $\displaystyle \lim _{n\to +\infty}\frac{1}{n^2}=0$. Donc $\displaystyle \lim _{n\to +\infty}u_n=\frac{1+0}{2-0}=\frac{1}{2}$. De même, $\displaystyle v_n=\frac{n\left(2+\frac{1}{n}\right)}{n^2\left(1-\frac{1}{n^2}\right)}=\frac{1}{n}\times\frac{2+\frac{1}{n}}{1-\frac{1}{n^2}}$ et donc $\displaystyle \lim _{n\to +\infty}v_n=0$.
\end{explanations}
\end{question}
}
